% GENERATED FILE. Edit canonical lessons, then run npm run generate:books.
% canonical-source-hash: fnv1a64:064cc893fca8b8e6
% canonical-lessons: ML-C456-ulppetuttuka, ML-C456-ulkkolluka, ML-C456-vertirikkuka, ML-C456-taramtirikkuka, ML-C456-vibhajikkuka

\chapter{Sorting Things Out}
\label{ch:ml-sorting-things-out}
\hlchaptermodality{456}

\begin{chapteropening}
\textbf{By the end of this chapter:} \emph{I can say to include, to hold, to separate, to sort and to divide.}

Complete the last lesson of chapter 456: I can say to include, to hold, to separate, to sort and to divide.
\end{chapteropening}

\section[\texorpdfstring{\ml{ഉൾപ്പെടുത്തുക} (uḷppeṭuttuka)}{ഉൾപ്പെടുത്തുക (uḷppeṭuttuka)}]{\ml{ഉൾപ്പെടുത്തുക} (uḷppeṭuttuka) — to include}

\label{lesson:ML-C456-ulppetuttuka}



\begin{quote}
Before the new one: say the Malayalam for ethical, then the Malayalam for to harm.
\end{quote}

\subsection*{You'll want to know: \ml{ഉൾപ്പെടുത്തുക}}
\textbf{\ml{ഉൾപ്പെടുത്തുക}} — \emph{uḷppeṭuttuka} — \textquotedblleft{}to include\textquotedblright{}.

\begin{quote}
\ml{അവർ} \ml{എന്റെ} \ml{പേരും} \ml{പട്ടികയിൽ} \ml{ഉൾപ്പെടുത്തി}.
\end{quote}

\begin{quote}
\emph{avar enṟe pēruṁ paṭṭikayil uḷppeṭutti.}
\end{quote}

\begin{quote}
They included my name in the list too.
\end{quote}

Past: \ml{ഉൾപ്പെടുത്തി}, \emph{uḷppeṭutti}.

The intransitive is \ml{ഉൾപ്പെടുക}, to be included.

\begin{cousinweb}
\textbf{\ml{ഉൾപ്പെടുത്തുക}} (\emph{uḷppeṭuttuka}) is a native Malayalam compound: \ml{ഉൾ}, \textquotedblleft{}inside\textquotedblright{}, and -\ml{പ്പെടുത്തുക}, \textquotedblleft{}to cause to be\textquotedblright{}, from \ml{പെടുക} (DEDR 3853): to cause to be inside.
\end{cousinweb}

\begin{grammarlens}[title={What you've built}]
One more word for saying what happens if.
\end{grammarlens}

\subsection*{Your turn}
\begin{itemize}
\raggedright
  \item \emph{Say it:} \emph{uḷppeṭuttuka}
  \item \emph{Say it:} \emph{uḷppeṭuttuka}, once more
  \item \emph{Say it:} \emph{avar enṟe pēruṁ paṭṭikayil uḷppeṭutti.} — They included my name in the list too.
  \item \emph{Recall:} say the Malayalam for ethical, then the Malayalam for to harm, then say \emph{uḷppeṭuttuka} again
\end{itemize}

\begin{tcolorbox}[breakable,colback=teal!4,colframe=teal!35!black,title={Before you move on}]
What does \ml{ഉൾപ്പെടുത്തുക} mean? (\textquotedblleft{}To include\textquotedblright{}.) Say it once more.
\end{tcolorbox}
\section[\texorpdfstring{\ml{ഉൾക്കൊള്ളുക} (uḷkkoḷḷuka)}{ഉൾക്കൊള്ളുക (uḷkkoḷḷuka)}]{\ml{ഉൾക്കൊള്ളുക} (uḷkkoḷḷuka) — to hold; to take in}

\label{lesson:ML-C456-ulkkolluka}



\begin{quote}
Before the new one: say the Malayalam for to harm, then the Malayalam for to include.
\end{quote}

\subsection*{You'll want to know: \ml{ഉൾക്കൊള്ളുക}}
\textbf{\ml{ഉൾക്കൊള്ളുക}} — \emph{uḷkkoḷḷuka} — \textquotedblleft{}to hold; to take in\textquotedblright{}.

\begin{quote}
\ml{ആ} \ml{ഹാളിൽ} \ml{അഞ്ഞൂറ്} \ml{പേരെ} \ml{ഉൾക്കൊള്ളാം}.
\end{quote}

\begin{quote}
\emph{ā hāḷil aññūṟŭ pēre uḷkkoḷḷāṁ.}
\end{quote}

\begin{quote}
That hall can hold five hundred people.
\end{quote}

Past: \ml{ഉൾക്കൊണ്ടു}, \emph{uḷkkoṇṭu}.

Also to accept an idea: \ml{ആ} \ml{ആശയം} \ml{അവൻ} \ml{ഉൾക്കൊണ്ടു}, he took that idea to heart.

\begin{cousinweb}
\textbf{\ml{ഉൾക്കൊള്ളുക}} (\emph{uḷkkoḷḷuka}) is a native Malayalam compound: \ml{ഉൾ}, \textquotedblleft{}inside\textquotedblright{}, and \ml{കൊള്ളുക}, \textquotedblleft{}to hold, to contain, to receive\textquotedblright{} (DEDR 2151): to take inside.
\end{cousinweb}

\begin{grammarlens}[title={What you've built}]
One more word for saying what happens if.
\end{grammarlens}

\subsection*{Your turn}
\begin{itemize}
\raggedright
  \item \emph{Say it:} \emph{uḷkkoḷḷuka}
  \item \emph{Say it:} \emph{uḷkkoḷḷuka}, once more
  \item \emph{Say it:} \emph{ā hāḷil aññūṟŭ pēre uḷkkoḷḷāṁ.} — That hall can hold five hundred people.
  \item \emph{Recall:} say the Malayalam for to harm, then the Malayalam for to include, then say \emph{uḷkkoḷḷuka} again
\end{itemize}

\begin{tcolorbox}[breakable,colback=teal!4,colframe=teal!35!black,title={Before you move on}]
What does \ml{ഉൾക്കൊള്ളുക} mean? (\textquotedblleft{}To hold; to take in\textquotedblright{}.) Say it once more.
\end{tcolorbox}
\section[\texorpdfstring{\ml{വേർതിരിക്കുക} (vērtirikkuka)}{വേർതിരിക്കുക (vērtirikkuka)}]{\ml{വേർതിരിക്കുക} (vērtirikkuka) — to separate; to tell apart}

\label{lesson:ML-C456-vertirikkuka}



\begin{quote}
Before the new one: say the Malayalam for to include, then the Malayalam for to hold.
\end{quote}

\subsection*{You'll want to know: \ml{വേർതിരിക്കുക}}
\textbf{\ml{വേർതിരിക്കുക}} — \emph{vērtirikkuka} — \textquotedblleft{}to separate; to tell apart\textquotedblright{}.

\begin{quote}
\ml{അവൾ} \ml{നല്ല} \ml{പഴങ്ങൾ} \ml{ചീത്തയിൽ} \ml{നിന്ന്} \ml{വേർതിരിച്ചു}.
\end{quote}

\begin{quote}
\emph{avaḷ nalla paḻaṅṅaḷ cīttayil ninnŭ vērtiriccu.}
\end{quote}

\begin{quote}
She separated the good fruit from the bad.
\end{quote}

Past: \ml{വേർതിരിച്ചു}, \emph{vērtiriccu}.

From something: -\ml{ൽ} \ml{നിന്ന്} \ml{വേർതിരിക്കുക}.

\begin{cousinweb}
\textbf{\ml{വേർതിരിക്കുക}} (\emph{vērtirikkuka}) is a native Malayalam compound: \ml{വേറ്}, \textquotedblleft{}separation, different\textquotedblright{} (DEDR 5548), and \ml{തിരിക്കുക}, \textquotedblleft{}to turn\textquotedblright{} (DEDR 3246): to turn things apart.
\end{cousinweb}

\begin{grammarlens}[title={What you've built}]
One more word for saying what happens if.
\end{grammarlens}

\subsection*{Your turn}
\begin{itemize}
\raggedright
  \item \emph{Say it:} \emph{vērtirikkuka}
  \item \emph{Say it:} \emph{vērtirikkuka}, once more
  \item \emph{Say it:} \emph{avaḷ nalla paḻaṅṅaḷ cīttayil ninnŭ vērtiriccu.} — She separated the good fruit from the bad.
  \item \emph{Recall:} say the Malayalam for to include, then the Malayalam for to hold, then say \emph{vērtirikkuka} again
\end{itemize}

\begin{tcolorbox}[breakable,colback=teal!4,colframe=teal!35!black,title={Before you move on}]
What does \ml{വേർതിരിക്കുക} mean? (\textquotedblleft{}To separate; to tell apart\textquotedblright{}.) Say it once more.
\end{tcolorbox}
\section[\texorpdfstring{\ml{തരംതിരിക്കുക} (taraṁtirikkuka)}{തരംതിരിക്കുക (taraṁtirikkuka)}]{\ml{തരംതിരിക്കുക} (taraṁtirikkuka) — to sort; to classify}

\label{lesson:ML-C456-taramtirikkuka}



\begin{quote}
Before the new one: say the Malayalam for to hold, then the Malayalam for to separate.
\end{quote}

\subsection*{You'll want to know: \ml{തരംതിരിക്കുക}}
\textbf{\ml{തരംതിരിക്കുക}} — \emph{taraṁtirikkuka} — \textquotedblleft{}to sort; to classify\textquotedblright{}.

\begin{quote}
\ml{അവൻ} \ml{പുസ്തകങ്ങൾ} \ml{വിഷയം} \ml{അനുസരിച്ച്} \ml{തരംതിരിച്ചു}.
\end{quote}

\begin{quote}
\emph{avan pustakaṅṅaḷ viṣayaṁ anusariccŭ taraṁtiriccu.}
\end{quote}

\begin{quote}
He sorted the books by subject.
\end{quote}

Past: \ml{തരംതിരിച്ചു}, \emph{taraṁtiriccu}.

\ml{തരം}, a kind, with \ml{തിരിക്കുക}, to turn: to turn things into kinds.

\begin{cousinweb}
\textbf{\ml{തരംതിരിക്കുക}} (\emph{taraṁtirikkuka}) is a Malayalam compound: \ml{തരം}, \textquotedblleft{}a kind, a sort\textquotedblright{}, whose origin is not settled, and \ml{തിരിക്കുക}, \textquotedblleft{}to turn\textquotedblright{}, native Dravidian (DEDR 3246).
\end{cousinweb}

\begin{grammarlens}[title={What you've built}]
One more word for saying what happens if.
\end{grammarlens}

\subsection*{Your turn}
\begin{itemize}
\raggedright
  \item \emph{Say it:} \emph{taraṁtirikkuka}
  \item \emph{Say it:} \emph{taraṁtirikkuka}, once more
  \item \emph{Say it:} \emph{avan pustakaṅṅaḷ viṣayaṁ anusariccŭ taraṁtiriccu.} — He sorted the books by subject.
  \item \emph{Recall:} say the Malayalam for to hold, then the Malayalam for to separate, then say \emph{taraṁtirikkuka} again
\end{itemize}

\begin{tcolorbox}[breakable,colback=teal!4,colframe=teal!35!black,title={Before you move on}]
What does \ml{തരംതിരിക്കുക} mean? (\textquotedblleft{}To sort; to classify\textquotedblright{}.) Say it once more.
\end{tcolorbox}
\section[\texorpdfstring{\ml{വിഭജിക്കുക} (vibhajikkuka)}{വിഭജിക്കുക (vibhajikkuka)}]{\ml{വിഭജിക്കുക} (vibhajikkuka) — to divide; to split}

\label{lesson:ML-C456-vibhajikkuka}



\begin{quote}
Before the new one: say the Malayalam for to separate, then the Malayalam for to sort.
\end{quote}

\subsection*{You'll want to know: \ml{വിഭജിക്കുക}}
\textbf{\ml{വിഭജിക്കുക}} — \emph{vibhajikkuka} — \textquotedblleft{}to divide; to split\textquotedblright{}.

\begin{quote}
\ml{അച്ഛൻ} \ml{സ്ഥലം} \ml{മൂന്ന്} \ml{മക്കൾക്കായി} \ml{വിഭജിച്ചു}.
\end{quote}

\begin{quote}
\emph{acchan sthalaṁ mūnnŭ makkaḷkkāyi vibhajiccu.}
\end{quote}

\begin{quote}
Father divided the land among his three children.
\end{quote}

Past: \ml{വിഭജിച്ചു}, \emph{vibhajiccu}.

The noun is \ml{വിഭജനം}, a division, a partition.

\begin{cousinweb}
\textbf{\ml{വിഭജിക്കുക}} (\emph{vibhajikkuka}) is built from Sanskrit \emph{vibhajana}, \textquotedblleft{}separation, distinction\textquotedblright{} (Monier-Williams), with the Malayalam verb ending -\ml{ിക്കുക}.
\end{cousinweb}

\begin{grammarlens}[title={What you've built}]
One more word for saying what happens if.
\end{grammarlens}

\subsection*{Your turn}
\begin{itemize}
\raggedright
  \item \emph{Say it:} \emph{vibhajikkuka}
  \item \emph{Say it:} \emph{vibhajikkuka}, once more
  \item \emph{Say it:} \emph{acchan sthalaṁ mūnnŭ makkaḷkkāyi vibhajiccu.} — Father divided the land among his three children.
  \item \emph{Recall:} say the Malayalam for to separate, then the Malayalam for to sort, then say \emph{vibhajikkuka} again
\end{itemize}

\begin{tcolorbox}[breakable,colback=teal!4,colframe=teal!35!black,title={Before you move on}]
What does \ml{വിഭജിക്കുക} mean? (\textquotedblleft{}To divide; to split\textquotedblright{}.) Say it once more.
\end{tcolorbox}
